\begin{problem}[2.]
    $z = 5 + 5 i \sqrt{3}$
\end{problem}

\begin{solution}
    Please refer to problem 1 for more detailed logic. 
    
    \[z = a + bi\]
    \[a = 5 \; \text{and} \; b = 5\sqrt{3} \]
    \[ \boxed{5 = Re(z)} \; \text{and} \; \boxed{ 5\sqrt{3} = Im(z) }\]
    
    Calculating r:
    \begin{align}
        r^2 &= a^2 + b^2 \\
        r^2 &= 5^2 + (5\sqrt{3})^2 \\
        r &= \sqrt{25 + 75} \\
        r &= \pm 10
    \end{align}
    
    \[\boxed{ |z| = 10 }\]
    
    According to \textbf{theorem 3}: Since $Re(z) \neq 0$ that means that $\tan(\theta) = \frac{Im(z)}{Re(z)} = \frac{b}{a}$.
    
    Calculating $\theta$:
    \begin{align}
        \tan(\theta) &= \frac{b}{a} \\ 
        \tan(\theta) &= \frac{5\sqrt{3}}{5} \\
        \tan(\theta) &= \sqrt{3}
    \end{align}
    \[\boxed{ \theta = \frac{\pi}{3} + 2 \pi k \; \text{for integers k} }\]
    
    \[\boxed{ arg(z) = \left\{ \frac{\pi}{3} + 2\pi k \; | \; \text{k is an integer} \right\} }\]
    
    \[\boxed{Arg(z) = \frac{\pi}{3}}\]
    
    \[\boxed{ \text{Polar Rep.} \; = 10 \; cis \left( \frac{\pi}{3} \right) }\]
    
\end{solution}